%% =====================================================================
\section{The group, and what it looks like $p$-adically}\label{sec:group}
%% =====================================================================

\subsection{The group, explicitly}\label{ss:groupexp}

We use Zudilin's realisation \cite[Lemma 8]{Zudilin2004} of the Rhin--Viola group for
weight $3$. Directions are pairs $(\alpha,\beta)=(\alpha_1,\dots,\alpha_4;\beta_1,\dots,\beta_4)$
of integer vectors with
\begin{equation}\label{eq:dircond}
 \{\beta_1,\beta_2\}<\{\alpha_1,\dots,\alpha_4\}<\{\beta_3,\beta_4\},
 \qquad \textstyle\sum_j\alpha_j=\sum_k\beta_k ,
\end{equation}
and the group acts by permutations of the sixteen entries of the matrix
\begin{equation}\label{eq:cmatrix}
 c_{jk}=|\alpha_j-\beta_k|,\qquad j,k=1,\dots,4 .
\end{equation}
%% REFEREE [R12]: Zudilin's Lemma 8 defines c_{jk} = a_j - b_k if a_j >= b_k and
%% b_k - a_j - 1 otherwise; under \eqref{eq:dircond} the two conventions differ by 1 in
%% exactly the entries c_33, c_44, both of which sit in Pi(c).  Harmless for delta by
%% Prop.~\ref{prop:deltainsens}, but the difference must be stated, not silently absorbed.
Zudilin's Lemma 8 writes $c_{jk}=b_k-a_j-1$ in the case $a_j<b_k$; under
\eqref{eq:dircond} that differs from \eqref{eq:cmatrix} by $1$ in the two entries
$c_{33},c_{44}$. Since the entries are scaled by $n$, the two conventions differ by an
$O(1)$ shift of the arguments of $\Pi$, which by Proposition~\ref{prop:deltainsens} leaves
$\dgr$ unchanged; we use \eqref{eq:cmatrix} throughout.

\begin{proposition}\label{prop:group}
\VER{BFS closure, exact} The group
$\Gp=\langle\mathfrak a_1,\mathfrak a_2,\mathfrak a_3,\mathfrak b,\mathfrak h\rangle$
generated by the involutions
\[
\begin{array}{ll}
 \mathfrak a_j\ (j=1,2,3): & \text{swap rows }j\text{ and }4,\\
 \mathfrak b: & \text{swap columns }3\text{ and }4,\\
 \mathfrak h: & (c_{11}\,c_{33})(c_{13}\,c_{31})(c_{22}\,c_{44})(c_{24}\,c_{42}),
\end{array}
\]
has order $1920=|W(D_5)|$.
\end{proposition}

The order is computed by breadth-first closure on the sixteen-entry configuration and agrees
with Zudilin's count. The group belongs to the Rhin--Viola series $W(A_4)$ of order $120$
for $\zeta(2)$, $W(D_5)$ of order $1920$ for $\zeta(3)$, $W(E_6)$ of order $51840$ for
$\zeta(4)$, and $\mathfrak S_7$ of order $5040$ for the Brown--Zudilin $\zeta(5)$ cell
integral \cite{BrownZudilin2022}.

For the hypergeometric parametrisation we pass to the variables of \cite[\S4]{Zudilin2004}.
With the normalisation $b_1=1$, set
\begin{equation}\label{eq:hvars}
 h_0=b_3+b_4-b_1-a_1,\qquad h_{1,2,3}=1-b_1+a_{2,3,4},\qquad
 h_4=b_4-a_1,\qquad h_5=b_3-a_1 .
\end{equation}
In these variables the generators act by
\begin{equation}\label{eq:hact}
\begin{aligned}
 \mathfrak a_j &:\ a_j\leftrightarrow a_4 \quad\text{(this moves }h_0\text{)},\\
 \mathfrak b &:\ b_3\leftrightarrow b_4,\\
 \mathfrak h &:\ (a_1,a_2,a_3,a_4;1,b_2,b_3,b_4)\\
 &\qquad\mapsto
   (b_3-a_3,\,a_2,\,b_3-a_1,\,a_4;\,1,\,b_2+b_3-a_1-a_3,\,b_3,\,b_3+b_4-a_1-a_3).
\end{aligned}
\end{equation}

\begin{remark}[why the first search fails]\label{rem:S5trivial}
The symmetric group $\mathfrak S_5$ permuting $h_1,\dots,h_5$ acts \emph{trivially} on the
rational function, because the function is a product over $j=1,\dots,5$. Hence the group is
invisible at fixed $h_0$. The nontrivial elements are the row swaps $\mathfrak a_j$, and
they \emph{move $h_0$}. Any search for transfer identities inside this family must be
cross-$h_0$; a search at fixed $h_0$ returns nothing and returns it for a trivial reason.
This observation is worth recording because it is the reason a natural first search fails,
and it also explains part of Finding~\ref{find:F3}.
\end{remark}

\subsection{The transfer identity, verified coefficientwise}\label{ss:transferver}

Combining \cite[eq.~(4.4) and Lemma 9]{Zudilin2004} gives the invariant in $h$-language. We
have not seen it stated in this form, and derive it here:
\begin{equation}\label{eq:invariant}
 \widetilde I:=\frac{H(c)}{\Pi(c)}
  =\frac{\widetilde G(a,b)}{\prod_j(a_j-b_1)!\cdot\prod_j(a_j-b_2)!}
  =\frac{\widetilde F(h)}{\prod_{j=1}^{5}(h_j-1)!\cdot\bigl(1+2h_0-\sum_jh_j\bigr)!},
\end{equation}
where
\begin{equation}\label{eq:Ftilde}
 \widetilde F(h)=\sum_{t\ge0}\widetilde R(h;t),\qquad
 \widetilde R(h;t)=(h_0+2t)\,(t+1)_{h_0-1}\prod_{j=1}^{5}\frac{1}{(t+h_j)_{1+h_0-2h_j}} .
\end{equation}
By Theorem~\ref{thm:transfer}, $\widetilde F(h)=A\zeta(3)+B$ with $A,B\in\QQ$ computed
exactly from the partial fractions of $\widetilde R$.

\begin{proposition}[anchor]\label{prop:aperyanchor}
\VER{$n\le4$, exact} At the Ball/Ap\'ery point $h_0=3n+2$, $h_j=n+1$ one has
$A/(2a_n)=1,\tfrac14,\tfrac1{36},\tfrac1{576}=1/n!^2$ for $n=1,2,3,4$, where
$a_n=\sum_k\binom nk^2\binom{n+k}k^2$ is Ap\'ery's numerator sequence \cite{Apery1979}.
\end{proposition}

So the machinery reproduces Ap\'ery's sequence exactly, confirming the
$2\,n!^2\sum\cdots$ normalisation of \cite{BallRivoal2001}.

\begin{finding}[the transfer identity]\label{find:transfer}
\VER{4 directions; orbits of sizes $1,40,80,1920$; $1321$ tested points; $0$ failures}
For every $g\in\Gp$ the invariant \eqref{eq:invariant} satisfies
$\widetilde I(h)=\widetilde I(gh)$ \emph{coefficientwise}, i.e.\ separately for $A/N$ and
for $B/N$, where $N$ is the factorial normaliser of \eqref{eq:invariant}. The verification:
\begin{center}
\tabsmall
\begin{tabular}{lccccc}
\toprule
direction & $|\Gp(a,b)|$ & $h$-valid points & identity holds & fails & distinct $\kappa$\\
\midrule
$(1,1,1,1;0,0,2,2)$ (symmetric) & $1$ & $1$ & $1$ & $0$ & $1$\\
$(2,2,3,3;0,1,4,5)$ & $40$ & $40$ & $40$ & $0$ & $3$\\
$(3,4,4,5;0,2,7,7)$ & $80$ & $80$ & $80$ & $0$ & $5$\\
$(4,7,8,11;0,3,13,14)$ & $1920$ & $1200$ & $1200$ & $0$ & $10$\\
\bottomrule
\end{tabular}
\end{center}
The $720$ skipped points at the last direction are those at which a brick degenerates,
$1+h_0-2h_j<1$; this is a direction-dependent boundary effect of the $h$-parametrisation,
not a failure of the identity.
\end{finding}

The ratio $\kappa=N(gh)/N(h)=\Pi(gc)/\Pi(c)$ is by construction a ratio of factorials. The
identity itself is Bailey's/Whipple's transformation for very-well-poised ${}_7F_6$,
\PROVED{} in \cite[Lemmas 7--9]{Zudilin2004}; what Finding~\ref{find:transfer} adds is the
exact coefficientwise check in the present normalisation.

\begin{finding}[half-integer robustness --- the $p$-adic point]\label{find:halfint}
\VER{3 directions; $780$ comparable orbit points; $0$ failures} Repeating the test with
$h_4,h_5\in\ZZ+\tfrac12$ --- testing proportionality of the full spectral vector by exact
$2\times2$ determinants over every pole-residue class --- gives
\begin{center}
\tabsmall
\begin{tabular}{lcc}
\toprule
direction & comparable orbit points & proportional / not\\
\midrule
$(2,2,3,3;0,1,4,5)$ & $120$ & $120$ / $0$\\
$(3,4,4,5;0,2,7,7)$ & $180$ & $180$ / $0$\\
$(4,7,8,11;0,3,13,14)$ & $480$ & $480$ / $0$\\
\bottomrule
\end{tabular}
\end{center}
\emph{The mechanism is lattice-robust: it is not an accident of integrality.}
\end{finding}

Finding~\ref{find:halfint} is what makes the Transfer Principle usable, because
$p$-adic admissibility forces half-integer parameters at $p=2$
(Proposition~\ref{prop:F1}). Without it, Corollary~\ref{cor:transfer} would be a statement
about a lattice the $p$-adic construction never visits.

\subsection{The savings function}\label{ss:savings}

\begin{definition}\label{def:delta}
With $F$ the eight-element index set of
$\Pi(c)=c_{21}!\,c_{31}!\,c_{41}!\,c_{12}!\,c_{32}!\,c_{42}!\,c_{33}!\,c_{44}!$, put
\begin{equation}\label{eq:deltadef}
\begin{aligned}
 \nu_\ell&=\max_{g\in\Gp}\ord_\ell\bigl(\Pi(cn)/\Pi(gcn)\bigr),\qquad
 \Phi_n=\prod_{\sqrt{m_0n}<\ell\le m_3n}\ell^{\nu_\ell},\\
 \Lambda(x)&=\max_{g\in\Gp}\sum_{i\in F}\bigl(\lfloor c_ix\rfloor-\lfloor(gc)_ix\rfloor\bigr),
 \qquad
 \dgr=\int_0^1\Lambda(x)\,\dd\psi(x)-\int_0^{1/m_3}\Lambda(x)\,\frac{\dd x}{x^2},
\end{aligned}
\end{equation}
$\psi$ the digamma function, and
\begin{equation}\label{eq:budget}
 \bud=2m_1+m_2,\quad m_1=\beta_4^*-\alpha_1^*,\quad
 m_2=\max\{\alpha_1-\beta_1,\ \alpha_2-\beta_2,\ \beta_4^*-\alpha_3,\ \beta_4^*-\alpha_4,\
   \beta_3^*-\alpha_1^*\}.
\end{equation}
\end{definition}

\begin{finding}[end-to-end validation of the savings pipeline]\label{find:rvvalidation}
\VER{$8$ digits, exact $\Lambda$, numerical quadrature of \eqref{eq:deltadef}} At Rhin and
Viola's optimum $(\alpha;\beta)=(18,17,16,19;\,0,7,31,32)$ the implementation returns
\[
\begin{array}{llll}
 m_0=19,& m_1=16,& m_2=18,& m_3=16,\\
 \bud=50, & \int_0^1\Lambda\,\dd\psi=24.18768530, & \int_0^{1/16}\Lambda\,\dd x/x^2=4.00000000, &\\
 \dgr=20.18768530, & C_2=29.81231470, & \text{orbit of distinct }F\text{-multisets}=118, &\\
 \multicolumn{4}{l}{\mu(\zeta(3))\le5.51389063\quad\text{against the published }5.51389062 .}
\end{array}
\]
\end{finding}

Eight-digit agreement with a published irrationality measure that was obtained by entirely
independent means is a strong joint test: the group, the invariant, the function $\Lambda$
and both integrals in \eqref{eq:deltadef} are simultaneously confirmed.

\subsection{The totally symmetric point banks nothing}\label{ss:symzero}

\begin{proposition}\label{prop:symzero}
\PROVED{} At the Ap\'ery/Ball point $\alpha=(1,1,1,1)$, $\beta=(0,0,2,2)$ every entry of the
matrix \eqref{eq:cmatrix} equals $|\alpha_j-\beta_k|=1$. A group acting by permutations of
sixteen equal entries has a single orbit point, so the stabiliser is all of $\Gp$ and
\[
 \nu_\ell\equiv0,\qquad \Phi_n\equiv1,\qquad \dgr=0\ \text{exactly}.
\]
The same holds at LSZ's $\zeta_2(5)$ point, where all bricks are again equal.
\end{proposition}

\begin{proof}
Every $g\in\Gp$ is a permutation of the sixteen entries $c_{jk}$; if all entries are equal
then $gc=c$ as a multiset and $\Pi(gcn)=\Pi(cn)$ for every $g$ and every $n$.
\end{proof}

Proposition~\ref{prop:symzero} is the good news and the bad news in one line.
\emph{The good news}: Ap\'ery, Beukers, Calegari and LSZ have banked exactly zero group
savings, so the orbit method is entirely fresh money. \emph{The bad news} is
\S\ref{ss:tension}: the locus where $\dgr=0$ is exactly the locus where the archimedean
coefficient growth $C_1$ also vanishes, and the $p$-adic criterion must pay $C_1$.

\subsection{How much money is there}\label{ss:howmuch}

%% REFEREE [R11]: the tag said "$\sum\alpha\le30$"; g_opt.py sweeps 9556 admissible
%% shapes with sum <= 34 and returns exactly the tabulated 0.412182 / 11.953292 / 29 /
%% 1.236547.  Range corrected upward to what the code actually does.
\begin{finding}\label{find:howmuch}
\VER{$9556$ integral directions with $\sum\alpha\le34$; hill-climbing to
$\sum\alpha\approx250$}
The scale-free relative saving $\dgr/\bud$, maximised over integral directions
satisfying \eqref{eq:dircond}:
\begin{center}
\tabsmall
\begin{tabular}{lcccc}
\toprule
directions & $\dgr$ & $\bud$ & $\dgr/\bud$ & $\dgr$ at a $\bud=3$ configuration\\
\midrule
$(4,7,8,11;0,3,13,14)$ & $11.9533$ & $29$ & $0.41218$ & $1.2365$\\
$(2,7,8,9;0,1,12,13)$ & $12.6730$ & $32$ & $0.39603$ & $1.1881$\\
Rhin--Viola optimum & $20.1877$ & $50$ & $0.40375$ & $1.2113$\\
symmetric & $0$ & $3$ & $0$ & $0$\\
\bottomrule
\end{tabular}
\end{center}
The absolute ceiling found, subject to the construction still working, is
$\dgr/\bud=0.4378$ at $(5,7,9,5;0,2,11,13)$.
\end{finding}

Consistency with the literature: Brown and Zudilin's $\zeta(5)$ record point has
$\dgr=34.394$ against $\sum m=84$, i.e.\ $0.4094$; Rhin and Viola's $\zeta(3)$ optimum gives
$0.4038$. Both land near $40\%$, and our independent maximisation caps out near $44\%$. It
is fair to summarise: \emph{a Rhin--Viola group removes about $40\%$ of the denominator,
fairly universally} \VER{the three directions above plus the hill-climb}.

\subsection{Four honest flags}\label{ss:flags}

The Transfer Principle transfers identities. It does not transfer admissibility, and it does
not manufacture a group. Here are the four places where the $p$-adic setting differs, stated
as they are, positive and negative.

\begin{proposition}[(F1) very-well-poisedness forces $\theta=\tfrac12$, i.e.\ $p=2$]\label{prop:F1}
\PROVED{} For $\Gamma(h_j+t)/\Gamma(1+h_0-h_j+t)$ to be a \emph{rational} function of $t$
one needs $2h_j-h_0-1\in\ZZ$. With $h_0\in\ZZ$ this forces $h_j\in\tfrac12\ZZ$. Hence the
only very-well-poised $p$-adic shift is $\theta=\tfrac12$. For $p\ge3$ one has
$|\theta|_p\ge q_p$, so $\theta=a/p^l$, and the shifts must come as the \emph{full} set
$\{\nu/p^l:p\nmid\nu\}$ --- which is exactly what \cite[Def.~16]{LaiSprang2023} and
\cite[\S7]{Lai2024} do, and it is why those $p$-adic constructions abandon the factor
$2t+h_0$.
\end{proposition}

A sanity check on Proposition~\ref{prop:F1}: Lai's normalising constant
$p^{(l+1/(p-1))M\varphi(p^l)n}$ of \cite[\S7]{Lai2024} collapses at $p=2$, $l=1$,
$\varphi(2)=1$, $M=4$ to $2^{8n}$ --- exactly LSZ's normalisation $C=2^{8n}$. The case
$p=2$ is cheap precisely because $\varphi(2)=1$.

\begin{proposition}[(F2) admissibility constrains the base point only]\label{prop:F2}
\PROVED{} For the Volkenborn integral at shift $\theta$ one needs
\begin{enumerate}[label=\textup{(R\arabic*)},leftmargin=2.8em]
\item all poles of $R$ at integers, else $R(t+\theta)$ has poles on $\Zp$;
\item all numerator bricks at shift $\equiv-\theta\pmod1$, else each numerator factor costs
      $-l$ instead of gaining $+1/(p-1)$ (see \eqref{eq:contrib} below).
\end{enumerate}
Orbit points need satisfy neither. They enter only through (i) the exact rational transfer
identity and (ii) their own denominator bound --- both statements about $\QQ$.
\end{proposition}

Proposition~\ref{prop:F2} is what makes the whole method conceivable, and it is the same
logic as the Brown--Zudilin saving
$\nu_p=\max_{g}\ord_p\prod_Fh_i!/\prod_F(gh_i)!$ \cite{BrownZudilin2022}.

\begin{finding}[(F3) the LSZ shape class has no internal transfer identity]\label{find:F3}
\EXCL{exact rational arithmetic, zero tolerance; cross-$h_0$; ranges below} Parametrise the
LSZ shape --- numerator bricks at $\theta$, denominator bricks at $\ZZ$, all centred, i.e.\
very-well-poised ---
\[
 R(t)=(2t+h_0)\prod_{j=1}^{M}(t+\theta+e_j)_{h_0-2e_j}\Big/
      \prod_{j=1}^{M}(t+f_j)_{h_0+1-2f_j}.
\]
An exhaustive exact search for proportional pairs of coefficient vectors --- the exact
signature of a transfer identity, by Corollary~\ref{cor:transfer} --- returns
\begin{center}
\tabsmall
\begin{tabular}{lcc}
\toprule
family & points tested & proportional pairs\\
\midrule
$M=2$ ($\zeta_p(3)$), all $h_0\le9$, cross-$h_0$ & $113$ & $0$\\
$M=4$ ($\zeta_p(5)$, LSZ), all $h_0\le7$, cross-$h_0$ & $1629$ & $0$\\
$M=4$, fixed $h_0=6$ & $672$ & $0$\\
%% REFEREE [R10]: the row was labelled "fixed $h_0\le6$".  g_search.py runs four
%% fixed-$h_0$ blocks, $h_0=3,4,5,6$, of sizes 15/127/127/672 (941 in all); 672 is the
%% $h_0=6$ block alone.  Relabelled, which also keeps the total 113+1629+672 = 2414 of
%% \S\ref{ss:summary} correct.  Both cross-$h_0$ rows reproduce exactly.
\bottomrule
\end{tabular}
\end{center}
\emph{One cannot find the group by deforming LSZ inside its own shape class: its group is
trivial.} This is a genuine negative, and it explains why the literature has not stumbled on
the group.
\end{finding}

\begin{finding}[(F4) genuinely $p$-adic points do have transfer partners]\label{find:F4}
\VER{full very-well-poised half-integer family; $5507$ distinct reduced rational functions,
$598$ genuinely $p$-adic; $13$ proportionality classes} Widen the search to the full
very-well-poised half-integer family ($h_0\in\ZZ$, $h_j\in\tfrac12\ZZ$, any $q$, reduced
modulo brick cancellation). Of the $598$ genuinely $p$-adic points --- integer poles
\emph{and} a non-empty half-integer numerator --- exactly $13$ lie in a proportionality class
containing a distinct partner. Four instances:
\[
\begin{array}{llll}
 h_0=5,\ \mathrm{num}\{\tfrac52,\tfrac52\},\ \mathrm{den}\{1^2,2^4,3^4,4^2\}
   &\leftrightarrow& h_0=4,\ \mathrm{den}\{1^4,2^4,3^4\} & \kappa=2\\
 h_0=6,\ \mathrm{num}\{\tfrac52,\tfrac72\},\ \mathrm{den}\{2,3^4,4\}
   &\leftrightarrow& h_0=4,\ \mathrm{den}\{1^3,3^3\} & \kappa=1\\
 h_0=7,\ \mathrm{num}\{\tfrac72,\tfrac72\},\ \mathrm{den}\{1,2,3^4,4^4,5,6\}
   &\leftrightarrow& h_0=5,\ \mathrm{den}\{1,2^4,3^4,4\} & \kappa=27/2\\
 h_0=6,\ \mathrm{num}\{\tfrac52,\tfrac72\},\ \mathrm{den}\{1,2^3,3^4,4^3,5\}
   &\leftrightarrow& h_0=7,\ \mathrm{num}\{1,2,5,6\},\ \mathrm{den}\{3^4,4^4\} & \kappa=2^8
\end{array}
\]
In every instance the partner is an \emph{integer-parameter} point, i.e.\ a member of the
classical archimedean family where the group and the whole denominator theory live.
\end{finding}

The mechanism is a quadratic transformation --- Gauss duplication,
$(t+\tfrac12)_n(t+1)_n=2^{-2n}(2t+1)_{2n}$ --- the classical bridge between half-integer and
integer very-well-poised parameters. This is the shape of the bridge one wants:

\begin{openlemma}[the Gauss-duplication bridge]\label{ol:bridge}
\OPENC{} The $\theta=\tfrac12$ very-well-poised family maps, by Gauss duplication, onto the
integer very-well-poised family, with an explicit factorial ratio $\kappa$; consequently a
$p$-adic linear form equals $\kappa$ times an integer-family linear form, and the integer
family carries the full $1920$-element orbit. \emph{Status:} \VER{the $13$ instances of
Finding~\ref{find:F4}, at small size}; the systematic family version is not proved here, and
it looks like a known quadratic transformation waiting to be quoted rather than a new
theorem.
\end{openlemma}

\subsection{Non-vanishing under orbit moves}\label{ss:nonvan}

$\Qp$ has no positivity, so non-vanishing of the linear form is a genuine constraint and
must be tracked. Under an orbit move it is free.

\begin{proposition}\label{prop:orbitnonvan}
\PROVED{} An orbit move multiplies the whole linear form by the nonzero rational
$\kappa=\Pi(gc)/\Pi(c)$. Hence the leading term is scaled and never cancelled;
$S(gh)=0\iff S(h)=0$, so an orbit move can neither create nor destroy vanishing; and the
input to Bel's criterion --- for LSZ, the Casoratian
$\rho_{n,0}\rho_{n+1,3}-\rho_{n+1,0}\rho_{n,3}=3\cdot2^{16n+18}/(n+1)^5$
\cite[Lemma 15(b)]{LSZ2025} --- is scaled by $\kappa_n\kappa_{n+1}\ne0$ and stays nonzero.
\end{proposition}

\begin{remark}\label{rem:bridgenonvan}
The (F4) bridges of Finding~\ref{find:F4} are \emph{not} orbit moves: the partners are
different rational functions. Non-vanishing must therefore be re-checked there. At every
instance computed, $\kappa\ne0$ (observed values
$1,2,4,16,\tfrac{27}2,18,54,\tfrac4{27},\tfrac43,2^5,2^8$), so no vanishing is introduced
\VER{the $13$ instances}.
\end{remark}

\subsection{Literature position: the empty intersection}\label{ss:corpus}

\begin{finding}\label{find:corpus}
\VER{cross-tabulation of a $37$-paper corpus} A full cross-tabulation of the corpus for
$p$-adic markers (\texttt{Volkenborn}, $\zeta_p$, $\Zp$) against group markers
(\texttt{Rhin}, \texttt{permutation group}, \texttt{group of order}) returns \emph{no paper
that uses a Rhin--Viola group $p$-adically}. The intersection is empty. The two files that
mention both state the matter as open: \cite[Final remarks]{LSZ2025} and
\cite[Rem.~47]{Brown2026}, both quoted in \S\ref{ss:questions}.
\end{finding}

The nearest existing object is \cite[\S\S7--8]{Lai2024}, the only genuinely multi-parameter
$p$-adic Volkenborn family, $V_n(t)$ with free
$(M;\delta_1,\dots,\delta_J;l;w_n)$. Its saving is of the form $\Phi_n^{-s/J}$ with
$\Phi_n$ an $\inf_y$ over floor functions --- a Chudnovsky--Rukhadze--Hata brick saving,
\emph{not} a $\max_{g\in\Gp}$ orbit saving. The two are independent and compose.

Finding~\ref{find:corpus} is the sense in which Theorem~\ref{thm:transfer} is new: not that
the underlying computation is deep, but that the intersection it inhabits was empty. The
scope of the finding is exactly the corpus swept; we make no claim about the literature
beyond it.
